\begin{abstract}
\noindent How attention improves perception has been debated for a quarter century:
proposals range from enhancement of the sensory signal, to reduction of external
noise and more efficient selection, to shifts in the observer's decision criterion
and reductions of spatial uncertainty. These accounts are not mutually exclusive ---
all are thought to contribute --- but behaviour alone cannot cleanly separate them.
Here we show that an affine-feedback recurrent vision-transformer architecture with a spatial
working-memory state, trained with reinforcement learning plus a temporal joint-embedding
predictive auxiliary loss (weight $0.5$), reproduces several signatures of primate visual attention across
separately trained checkpoints. Because its internal attention can be read out and perturbed,
the model directly operationalizes two components of attentional performance: its decision
criterion shifts toward liberal reporting and its perceptual sensitivity ($d'$) depends on
the attention clamp. Injected attention also evokes false detections preferentially at the
monitored location, demonstrating spatial gating. These measurements do not constitute a
quantitative three-way decomposition and include no dedicated external-noise test. The model
deploys attention to a cued location, uses spatial validity when validity is the operative cue,
maintains a value code from cue onset, and develops a decodable, position-specific code for the
change when the change occurs. In a preliminary cued analogue of a paradigm used in the primate
literature, separately trained reward-regime checkpoints reproduce the Luo--Maunsell
sensitivity/criterion contrast. This RViT+ study extends the 2025 empirical Recurrent ViT
\cite{morgan_albanna_herman2025_rvit}; it is also a companion to, but not a revision
of, a separate 2026 normative account of when value-directed attention matters
\cite{morgan_albanna_herman2026_vda}.
\end{abstract}
